% !TeX root = eleve.tex
% % !TeX root = corrige.tex

% ========================
% BARÈME (affiché élève & corrigé)
% ========================

\begin{center}
	\badge{primary}{Ex.1: 4pts}\quad
	\badge{primary}{Ex.2: 4pts}\quad
	\badge{primary}{Ex.3: 5pts}\quad
	\badge{primary}{Ex.4: 4pts}\quad
	\badge{primary}{Ex.5: 3pts}
\end{center}

% ========================
% CONSIGNES
% ========================
\vspace{0.2cm}
\begin{consignebox}
	\faIcon{info-circle} \textbf{CONSIGNES IMPORTANTES}
	\begin{itemize}[leftmargin=*,noitemsep,topsep=5pt]
		\item[\faIcon{check}] Écrivez lisiblement et justifiez vos réponses
		\item[\faIcon{check}] Respectez l'orthographe des nombres en lettres
		\item[\faIcon{check}] Les exercices peuvent être traités dans n'importe quel ordre
		%\item[\faIcon{times}] Calculatrice non autorisée
	\end{itemize}
\end{consignebox}

% ========================
% EXERCICE 1
% ========================
\vspace{0.5cm}
\begin{exercisebox}{\faIcon{sort-numeric-down} Exercice 1 : Rang des chiffres \points{4}}
	Considérons le nombre : \textbf{\num{7256394}}
	
	\begin{enumerate}[label=\textcolor{primary}{\textbf{\alph*)}},leftmargin=*]
		\item Quel est le chiffre des centaines ? \dotfill \rep{3} \points{1}
		\item Quel est le chiffre des dizaines de milliers ? \dotfill \rep{5} \points{1}
		\item Quel est le \emph{nombre de milliers} dans ce nombre ? \dotfill \rep{\num{7256}} \points{1}
		\item Donne la décomposition décimale de ce nombre : \points{1}
		\rep{$7\times\num{1000000} + 2\times\num{100000} + 5\times\num{10000} + 6\times\num{1000} + 3\times\num{100} + 9\times\num{10} + 4$}
		\zonereponseLignes[0.4cm]{2}
	\end{enumerate}
\end{exercisebox}

% ========================
% EXERCICE 2
% ========================
\begin{exercisebox}{\faIcon{spell-check} Exercice 2 : Écriture en lettres \points{4}}
	\textbf{Écris les nombres suivants en toutes lettres :}
	
	\begin{enumerate}[label=\textcolor{primary}{\textbf{\alph*)}},leftmargin=*]
		\item \num{1200} : \dotfill \rep{mille-deux-cents} \points{1}
		\item \num{5087} : \dotfill \rep{cinq-mille-quatre-vingt-sept} \par \dotfill \points{1}
		\item \num{14680} : \dotfill \rep{quatorze-mille-six-cent-quatre-vingts} \par \dotfill \points{1}
		\item \num{300092} : \dotfill \rep{trois-cent-mille-quatre-vingt-douze} \par \dotfill \points{1}
	\end{enumerate}
\end{exercisebox}

\clearpage
% ========================
% EXERCICE 3
% ========================
\begin{exercisebox}{\faIcon{ruler-horizontal} Exercice 3 : Demi-droite graduée \points{5}}
	\vspace{0.3cm}
	\textbf{Observe la demi-droite graduée ci-dessous :}
	
	\begin{center}
		\begin{tikzpicture}[scale=1.2]
			% Axe (commence à 1200, va jusqu'à 2400)
			\draw[thick,->] (0,0) -- (13,0);
			
			% Graduations principales tous les 200
			\foreach \x/\val in {0/1200, 2/1400, 4/1600, 6/1800, 8/2000, 10/2200, 12/2400} {
				\draw[thick] (\x,0.15) -- (\x,-0.15);
				\node[below] at (\x,-0.4) {\num{\val}};
			}
			
			% Graduations secondaires tous les 100
			\foreach \x in {1,3,5,7,9,11} {
				\draw (\x,0.08) -- (\x,-0.08);
			}
			
			% Graduations tertiaires tous les 50
			\foreach \x in {0.5,1.5,2.5,3.5,4.5,5.5,6.5,7.5,8.5,9.5,10.5,11.5,12.5} {
				\draw (\x,0.04) -- (\x,-0.04);
			}
			
			% Points à placer
			\node[circle,fill=primary,inner sep=2.5pt,label=above:$A$] at (1.5,0) {};
			\node[circle,fill=primary,inner sep=2.5pt,label=above:$B$] at (5.5,0) {};
			\node[circle,fill=primary,inner sep=2.5pt,label=above:$C$] at (8.25,0) {};
			\node[circle,fill=primary,inner sep=2.5pt,label=above:$D$] at (10.75,0) {};
			
			% Point E pour le corrigé (à 2125)
			\onlycorrige{\node[circle,fill=warning,inner sep=2.5pt,label=above:$E$] at (9.25,0) {};}
		\end{tikzpicture}
	\end{center}
	
	\begin{enumerate}[label=\textcolor{primary}{\textbf{\alph*)}},leftmargin=*]
		\item Détermine l'unité de longueur de cette demi-droite. \points{1}
		\rep{L'unité de longueur est de 200 (chaque graduation principale correspond à 200 unités et chaque petite graduation 50)}
		\zonereponseLignes[0.3cm]{1}
		
		\item Quelle est l'abscisse du point $A$ ? \dotfill \rep{\num{1350}} \points{1}
		
		\item Quelle est l'abscisse du point $B$ ? \dotfill \rep{\num{1750}} \points{1}
		
		\item Quelle est l'abscisse du point $C$ ? \dotfill \rep{\num{2025}} \points{1}
		
		\item Place le point $E$ d'abscisse \num{2125} sur la demi-droite ci-dessus. \rep{Point $E$ placé entre 2000 et 2200, plus proche de 2200} \points{1}
	\end{enumerate}
\end{exercisebox}

% ========================
% EXERCICE 4
% ========================
\begin{exercisebox}{\faIcon{greater-than-equal} Exercice 4 : Comparaison et rangement \points{4}}
	\begin{enumerate}[label=\textcolor{primary}{\textbf{\alph*)}},leftmargin=*]
		\item Compare avec $<$, $>$ ou $=$ : \points{2}
		\begin{multicols}{4}
			\begin{itemize}[label=$\bullet$]
				\item \num{3456} \rep{$<$} \dotfill \num{3546}
				\item \num{12089} \rep{$<$} \dotfill \num{12098}
				\item \num{7000} \rep{$>$} \dotfill \num{6999}
				\item \num{45230} \rep{$>$} \dotfill \num{45203}
			\end{itemize}
		\end{multicols}
		\item Range dans l'ordre croissant : \points{2}
		\begin{center}
			\num{8743}\,;\quad \num{8347}\,;\quad \num{8734}\,;\quad \num{8437}\,;\quad \num{8374}
		\end{center}
		\rep{\num{8347} < \num{8374} < \num{8437} < \num{8734} < \num{8743}}
		\zonereponseLignes[0.4cm]{2}
	\end{enumerate}
\end{exercisebox}

% ========================
% EXERCICE 5
% ========================
\begin{exercisebox}{\faIcon{lightbulb} Exercice 5 : Problème \points{3}}
	\begin{tcolorbox}[colback=blue!5,colframe=primary!30,arc=2mm,boxrule=0.5pt]
		\textbf{Contexte :} La bibliothèque du collège compte ses livres.
		\begin{itemize}[noitemsep,topsep=5pt]
			\item Il y a \textbf{3 milliers} de romans
			\item Il y a \textbf{15 centaines} de livres documentaires
			\item Il y a \textbf{85 dizaines} de bandes dessinées
		\end{itemize}
	\end{tcolorbox}
	
	\begin{enumerate}[leftmargin=*]
		\item \textbf{Calcule} le nombre total de livres : \points{2}
		\rep{Romans : \num{3000} ; Documentaires : \num{1500} ; BD : \num{850} \\
			Total : \num{3000} + \num{1500} + \num{850} = \num{5350} livres}
		\zonereponseLignes[0.4cm]{3}
		\item \textbf{Écris} ce nombre total en toutes lettres : \points{1}
		\rep{cinq-mille-trois-cent-cinquante}
		\zonereponseLignes[0.4cm]{2}
	\end{enumerate}
\end{exercisebox}

% ========================
% AUTO-ÉVALUATION
% ========================
\vspace{0.3cm}
\begin{center}
	\begin{tikzpicture}
		\node[draw=secondary,rounded corners=3mm,inner sep=8pt,fill=light!30] {
			\begin{tabular}{lccccc}
				\textbf{Auto-évaluation :} &
				\faIcon{frown} Difficile &
				\faIcon{meh} Moyen &
				\faIcon{smile} Facile &
				\phantom{xx} &
				\textbf{Temps utilisé :} \underline{\phantom{xxxxx}} min
			\end{tabular}
		};
	\end{tikzpicture}
\end{center}